% SPDX-License-Identifier: GPL-3.0-or-later
% Copyright (C) 2026 Alessandro Greco
%% Descripción alternativa accesible (EPUB).
\altfig{Esquema didáctico de algunas crates Arti: arriba el programa arti con proxy SOCKS y arti-client; en el centro gestores de directorios, circuitos, guardias y servicios onion; abajo protocolo, celdas y criptografía. Las flechas continuas representan relaciones funcionales simplificadas, no el grafo exacto de dependencias Cargo. Las líneas discontinuas comparan funciones con módulos educativos de torlab (directory.py, path\_selection.py, onion.py, channel.py, circuit.py, cells.py, ntor.py y encryption.py), sin implicar dependencia ni equivalencia completa.}
% Principales crates de Arti y módulos correspondientes de torlab
\begin{tikzpicture}[font=\sffamily\scriptsize,
  crate/.style={draw=torviola,fill=torviola!8,rounded corners=2pt,minimum width=2.5cm,minimum height=0.6cm},
  lab/.style={draw=torverde!60!black,fill=torverde!10,rounded corners=2pt,minimum width=2.2cm,minimum height=0.6cm,font=\ttfamily\scriptsize}]
  \node[crate] (arti) at (0,4.2) {arti (SOCKS, línea de órdenes)};
  \node[crate] (client) at (0,3.3) {arti-client};
  \node[crate] (dirmgr) at (-2.9,2.3) {tor-dirmgr};
  \node[crate] (circmgr) at (0,2.3) {tor-circmgr};
  \node[crate] (guardmgr) at (2.9,2.3) {tor-guardmgr};
  \node[crate] (netdoc) at (-2.9,1.3) {tor-netdoc};
  \node[crate] (chanmgr) at (0,1.3) {tor-chanmgr};
  \node[crate] (hs) at (2.9,1.3) {tor-hsclient};
  \node[crate] (proto) at (0,0.3) {tor-proto};
  \node[crate] (cell) at (0,-0.7) {tor-cell};
  \node[crate] (llc) at (0,-1.7) {tor-llcrypto};
  \draw[->] (arti) -- (client);
  \draw[->] (client) -- (dirmgr); \draw[->] (client) -- (circmgr); \draw[->] (client) -- (hs);
  \draw[->] (dirmgr) -- (netdoc); \draw[->] (circmgr) -- (chanmgr); \draw[->] (circmgr) -- (guardmgr);
  \draw[->] (chanmgr) -- (proto); \draw[->] (proto) -- (cell); \draw[->] (cell) -- (llc);
  \draw[->] (dirmgr) -- (circmgr); \draw[->] (hs) -- (circmgr);
  % Mapa conceptual: los trazos discontinuos son analogías, no dependencias Cargo.
  \node[lab] (d) at (-6.2,1.3) {directory.py};
  \node[lab] (p) at (6.0,2.3) {path\_selection.py};
  \node[lab] (o) at (6.0,1.3) {onion.py};
  \node[lab] (c) at (-4.4,0.3) {channel.py, circuit.py};
  \node[lab] (ce) at (-4.4,-0.7) {cells.py};
  \node[lab] (n) at (-4.4,-1.7) {ntor.py, encryption.py};
  \foreach \a/\b in {d/netdoc,p/guardmgr,o/hs,c/proto,ce/cell,n/llc} \draw[dashed,torgrigio] (\a) -- (\b);
  \node[torviola] at (0,5) {Arti (Rust)};
  \node[torverde!50!black] at (-5.2,2.4) {torlab (Python)};
\end{tikzpicture}
